% FORMAT AND PACKAGES
% {
\documentclass[a4paper]{article}
\usepackage{a4wide,amssymb,epsfig,latexsym,multicol,array,hhline,fancyhdr}
\usepackage{tcolorbox}
\usepackage{minted}
\usepackage{vntex}
\usepackage{amsmath}
\usepackage{lastpage}
\usepackage[lined,boxed,commentsnumbered]{algorithm2e}
\usepackage{enumerate}
\usepackage{xcolor}
\usepackage{graphicx}							% Standard graphics package
\usepackage{array}
\usepackage{tabularx, caption}
\usepackage{multirow}
\usepackage{multicol}
\usepackage{rotating}
\usepackage{graphics}
\usepackage{geometry}
\usepackage{setspace}
\usepackage{epsfig}
\usepackage{tikz}
\usepackage{xfrac}
\usepackage{bm}
\usepackage{biblatex}
\usepackage[colorlinks]{hyperref}
\newcommand{\cach}{\hspace*{1.5em}\ignorespaces}
% \usepackage[acronym,toc]{glossaries}
% \usepackage[symbols,nogroupskip,nonumberlist]{glossaries-extra}
\usepackage[
 sort=none,% no sorting or indexing required
 abbreviations,% create list of abbreviations
 symbols,% create list of symbols
 stylemods,style=list, % set the default glossary style
 nogroupskip, nonumberlist, nomain
]{glossaries-extra}


% FORMATTING
% {
\DeclareMathOperator{\arccot}{arccot}
\captionsetup[table]{name=Bảng}
\captionsetup[figure]{name=Hình}
\newenvironment{Description}{\list{}{%
    \let\makelabel\descriptionlabel    % this comes from the original description environment
    \setlength{\rightmargin}{\leftmargin}% this comes from the original quote environment
    \setlength{\labelwidth}{0pt}%          this is new
    }}{\endlist}

\addbibresource{citations.bib}
    
\hypersetup{urlcolor=blue,linkcolor=black,citecolor=black,colorlinks=true} 
\usetikzlibrary{arrows,snakes,backgrounds}
\definecolor{mathblue}{RGB}{0,114,188}
% \makeatletter  \def\m@th{\mathsurround\z@\color{mathblue}} \makeatother
% \everymath{\color{mathblue}}
% \setmathfont[Color=000000]{Arial}
%\usepackage{pstcol} 								% PSTricks with the standard color package
\newtheorem{theorem}{{\bf Theorem}}
\newtheorem{property}{{\bf Property}}
\newtheorem{proposition}{{\bf Proposition}}
\newtheorem{corollary}[proposition]{{\bf Corollary}}
\newtheorem{lemma}[proposition]{{\bf Lemma}}

\AtBeginDocument{\renewcommand{\listfigurename}{Danh sách hình ảnh}}
\AtBeginDocument{\renewcommand{\listtablename}{Danh sách bảng biểu}}
\AtBeginDocument{\renewcommand*\contentsname{Mục lục}}
\AtBeginDocument{\renewcommand*\refname{Tài liệu tham khảo}}
%\usepackage{fancyhdr}

\setlength{\headheight}{40pt}
\pagestyle{fancy}
\fancyhead{} % clear all header fields
\fancyhead[L]{
 \begin{tabular}{rl}
    \begin{picture}(25,15)(0,0)
    \put(0,-8){\includegraphics[width=8mm, height=8mm]{hcmut.png}}
    %\put(0,-8){\epsfig{width=10mm,figure=hcmut.eps}}
   \end{picture}&
	%\includegraphics[width=8mm, height=8mm]{hcmut.png} & %
	\begin{tabular}{l}
		\textbf{\bf \ttfamily Trường Đại học Bách Khoa TP.Hồ Chí Minh}\\
		\textbf{\bf \ttfamily Khoa Khoa học và Kỹ thuật máy tính}
	\end{tabular} 	
 \end{tabular}
}
\fancyhead[R]{
	\begin{tabular}{l}
		\tiny \bf \\
		\tiny \bf 
	\end{tabular}  }
\fancyfoot{} % clear all footer fields
\fancyfoot[L]{\scriptsize \ttfamily Báo cáo Bài tập lớn môn Mô hình hóa toán học - HK241 - Năm học 2024 - 2025}
\fancyfoot[R]{\scriptsize \ttfamily Trang {\thepage}/\pageref{LastPage}}
\renewcommand{\headrulewidth}{0.3pt}
\renewcommand{\footrulewidth}{0.3pt}

\setcounter{secnumdepth}{4}
\setcounter{tocdepth}{4}

\makeatletter
\newcounter {subsubsubsection}[subsubsection]
\renewcommand\thesubsubsubsection{\thesubsubsection .\@alph\c@subsubsubsection}
\newcommand\subsubsubsection{\@startsection{subsubsubsection}{4}{\z@}%
                                     {-3.25ex\@plus -1ex \@minus -.2ex}%
                                     {1.5ex \@plus .2ex}%
                                     {\normalfont\normalsize\bfseries}}
\newcommand*\l@subsubsubsection{\@dottedtocline{3}{10.0em}{4.1em}}
\newcommand*{\subsubsubsectionmark}[1]{}
% \def\m@th{\mathsurround\z@\color{mathblue}}
\makeatother
% }
% }

% ACRONYMS & SYMBOLS
% {
% \makeglossaries
\setabbreviationstyle{long-short}
\newabbreviation{csp}{CSP}{Cutting Stock Problem}
\newabbreviation{ffd}{FFD}{First Fit Decreasing}
\newabbreviation{ga}{GA}{Genetic Algorithm}
\newabbreviation{lp}{LP}{Linear Programming}
% \glsnoexpandfields
\glsxtrnewsymbol[description = {Tập hợp số tự nhiên}]{natural}{\ensuremath{\mathbb{N}}}

% }
%
% DOCUMENT
\begin{document}

% TITLE PAGE
\begin{titlepage}
\begin{center}
ĐẠI HỌC QUỐC GIA THÀNH PHỐ HỒ CHÍ MINH\\
TRƯỜNG ĐẠI HỌC BÁCH KHOA\\
KHOA KHOA HỌC VÀ KỸ THUẬT MÁY TÍNH\\
\end{center}

\vspace{1cm}

\begin{figure}[h!]
\begin{center}
\includegraphics[width=3cm]{hcmut.png}
\end{center}
\end{figure}

\vspace{1cm}


\begin{center}
\begin{tabular}{c}
\multicolumn{1}{c}{\textbf{{\Large Đồ án tổng hợp - CNPM (CO3103)}}}\\
~~\\
\hline
\\
\multicolumn{1}{l}{\textbf{{\Large Bài tập lớn}}}\\
\\
\textbf{\textit{{\Huge "Ứng dụng nghe nhạc trực tuyến"}}}\\
\\
\hline
\end{tabular}
\end{center}

\vspace{2cm}

\begin{table}[h]
\centering
    \begin{tabular}{rl}
    \hspace{3 cm}\textbf{GVHD}:
    & ThS. Trần Trương Tuấn Phát\\

    & \\[10pt]
\textbf{Sinh viên}: & Lư Chấn Vũ - 2313955 \emph{(Nhóm 10, \textbf{Leader})} \\
& Nguyễn Phú Vinh - 2313922 \emph{(Nhóm 10)} \\
& Trần Dương Khiết Nhi - 2312509 \emph{(Nhóm 10)} \\
& Lê Minh Khoa - 2311593 \emph{(Nhóm 10)} \\
& Lê Minh Trí - 2313593 \emph{(Nhóm 10)} \\
    \end{tabular}
\end{table}

\begin{center}
{\footnotesize TP. HỒ CHÍ MINH, 09/2025}
\end{center}
\end{titlepage}

\pagebreak
\tableofcontents

\pagebreak

% Glossaries
% {}
\printunsrtglossary[type={symbols}, title={Danh sách kí hiệu}]
\printunsrtglossary[type={abbreviations}, title={Danh sách từ viết tắt}]
\pagebreak
\listoffigures
\listoftables
\pagebreak
\addcontentsline{toc}{section}{\listfigurename}
\addcontentsline{toc}{section}{\listtablename}

% 

% Member list
\section*{Danh sách thành viên và nhiệm vụ}
\addcontentsline{toc}{section}{Danh sách thành viên và nhiệm vụ}
\begin{center}
\begin{table}[H]
\centering
\begin{tabular}{|c|c|c|l|c|}
\hline
\textbf{STT} & \textbf{Họ và tên} & \textbf{MSSV} & \textbf{Nhiệm vụ} & \textbf{\% hoàn thành}\\
\hline 
%%%%%Student 1%%%%%%%%%%
\multirow{3}{*}{1} & \multirow{3}{*}{Lư Chấn Vũ} & \multirow{3}{*}{2313955} & 
- & \multirow{3}{*}{100\%}\\
 & & & - & \\
\hline
%%%%%Student 2%%%%%%%%%%
\multirow{3}{*}{2} & \multirow{3}{*}{Nguyễn Phú Vinh} & \multirow{3}{*}{2313922} & 
- & \multirow{3}{*}{100\%}\\
 & & & - & \\
\hline
%%%%%Student 3%%%%%%%%%%
\multirow{3}{*}{3} & \multirow{3}{*}{Trần Dương Khiết Nhi} & \multirow{3}{*}{2312509} & 
-& \multirow{3}{*}{100\%}\\
 & & & - & \\
\hline
%%%%%Student 4%%%%%%%%%%
\multirow{3}{*}{4} & \multirow{3}{*}{Lê Minh Khoa} & \multirow{3}{*}{2311593} & 
- & \multirow{3}{*}{100\%}\\
 & & & - & \\
\hline
%%%%%Student 5%%%%%%%%%%
\multirow{3}{*}{5} & \multirow{3}{*}{Lê Minh Trí} & \multirow{3}{*}{2313593} & 
- & \multirow{3}{*}{100\%}\\
 & & & - & \\
\hline
\end{tabular}
\caption{\label{table1}Danh sách thành viên và nhiệm vụ}
\end{table}
\end{center}

\pagebreak
\section*{Nội dung báo cáo}

% Non-functional requirements
\begin{enumerate}
    % Hiệu năng (Performance)
    \item \textbf{Hiệu năng (Performance):} Hệ thống phải có khả năng phản hồi nhanh chóng trong mọi tình huống sử dụng phổ biến.
    \begin{itemize}
        \item Thời gian phản hồi cho các thao tác chính (như đăng nhập, tìm kiếm bài hát, phát nhạc, chuyển bài, tạo playlist,...) \textbf{không được vượt quá 2 giây} (đối với tối thiểu 95\% yêu cầu), ngay cả khi có \textbf{1000 người dùng đồng thời} truy cập.
        \item Thời gian bắt đầu phát nhạc (từ khi ấn nút Play đến khi nghe) \textbf{$\leq \mathbf{800}$ ms} trong điều kiện mạng ổn định.
        \item Tác vụ nhận diện bài hát từ đoạn ghi âm 6–10s phải trả kết quả trong vòng \textbf{$\leq \mathbf{2}$ giây} để đảm bảo trải nghiệm mượt mà.
    \end{itemize}

    % Bảo mật (Security)
    \item \textbf{Bảo mật (Security):} Hệ thống cần đảm bảo mức độ bảo mật cao.
    \begin{itemize}
        \item Người dùng phải đăng nhập qua tài khoản riêng (email/social login OAuth 2.0) để đảm bảo xác thực an toàn.
        \item Dữ liệu cá nhân (tên, email, playlist, lịch sử nghe) và dữ liệu nhạc phải được mã hóa khi lưu trữ (AES-256) và truyền tải (TLS 1.2+).
        \item Hệ thống áp dụng \textbf{role-based access control} (phân quyền): người dùng thường không thể truy cập chức năng quản trị.
        \item Phiên đăng nhập tự động hết hạn sau \textbf{30 phút không hoạt động}.
        \item Hệ thống có cơ chế giới hạn tốc độ (rate limit) để ngăn chặn brute force và tấn công DDoS:
        \begin{itemize}
            \item Tối đa 5 lần đăng nhập thất bại trong vòng 15 phút sẽ khóa tài khoản trong 30 phút.
            \item Giới hạn tối đa 100 yêu cầu API mỗi phút cho mỗi địa chỉ IP.
            \item Mỗi tài khoản người dùng chỉ được phép phát tối đa 3 luồng nhạc đồng thời.
            \item Giới hạn tối đa 10 yêu cầu nhận diện bài hát mỗi phút cho mỗi tài khoản người dùng.
        \end{itemize}
    \end{itemize}

    % Độ tin cậy (Reliability)
    \item \textbf{Độ tin cậy (Reliability):} Hệ thống phải có khả năng hoạt động liên tục, ổn định và không bị gián đoạn.
    \begin{itemize}
        \item Hệ thống cần đảm bảo \textbf{độ sẵn sàng (uptime)} tối thiểu là \textbf{99.9\%} cho dịch vụ phát nhạc, và tối thiểu \textbf{99.5\%} cho dịch vụ nhận diện bài hát.
        \item Cơ chế sao lưu dữ liệu được thực hiện tự động hàng ngày để tránh mất mát thông tin trong trường hợp xảy ra sự cố hệ thống.
        \item Có thông báo lỗi rõ ràng nếu hệ thống không khả dụng, kèm theo phương án xử lý.
        \item Khi gặp sự cố, phải có cơ chế \textbf{khôi phục dữ liệu trong vòng 24 giờ}.
    \end{itemize}

    % Khả năng sử dụng (Usability)
    \item \textbf{Khả năng sử dụng (Usability):}
    \begin{itemize}
        \item Giao diện người dùng (UI) trực quan, thân thiện, quen thuộc với các biểu tượng phổ biến (Play/Pause, Next, Like, Playlist,...).
        \item Hỗ trợ \textbf{song ngữ (tiếng Việt và tiếng Anh)} để phù hợp với nhiều đối tượng người dùng.
        \item Người dùng mới có thể làm quen và thành thạo với các chức năng cơ bản (tìm kiếm, phát nhạc, playlist, nhận diện bài hát,...) sau \textbf{$\leq \mathbf{15}$ phút} hướng dẫn.
        \item Hỗ trợ \textbf{đa nền tảng}: trình duyệt web phổ biến (Chrome, Firefox, Edge, Safari – 2 phiên bản gần nhất) và ứng dụng di động (Android 10+, iOS 14+).
        \item Hỗ trợ \textbf{accessibility} (screen reader, contrast mode, điều khiển bàn phím,...) để đảm bảo các đối tượng khuyết tật có thể sử dụng.
    \end{itemize}

    % Khả năng mở rộng (Scalability)
    \item \textbf{Khả năng mở rộng (Scalability):}
    \begin{itemize}
        \item Hệ thống có khả năng phục vụ \textbf{thêm ít nhất 3000 người dùng/năm} mà không ảnh hưởng hiệu năng.
        \item Có thể mở rộng, tích hợp thêm các module mới như: gợi ý nhạc bằng AI, podcast,... mà không làm thay đổi cấu trúc hiện tại.
        \item Các microservice (nghe nhạc, tìm kiếm, nhận diện, gợi ý) có thể scale độc lập trên cloud.
    \end{itemize}

    % Tính sẵn sàng (Availability)
    \item \textbf{Tính sẵn sàng (Availability):}
    \begin{itemize}
        \item Hệ thống cần hoạt động liên tục và vận hành ổn định \textbf{24/7} để hỗ trợ người dùng mọi lúc.
        \item Thời gian bảo trì định kỳ \textbf{không quá 4 lần/năm}, mỗi lần $\leq 4$ giờ.
        \item Nếu bảo trì, cần có thông báo trên hệ thống trước ít nhất \textbf{24 giờ}.
        \item Trong thời gian bảo trì, hệ thống phải hiển thị thông báo rõ ràng cho người dùng.
    \end{itemize}

    % Tính tương thích (Compatibility)
    \item \textbf{Tính tương thích (Compatibility):}
    \begin{itemize}
        \item Ứng dụng tương thích với nhiều thiết bị (điện thoại, máy tính, máy tính bảng) và nhiều chuẩn phát nhạc (HLS/DASH, AAC/Opus).
        \item Hệ thống hoạt động tốt trên nhiều nền tảng thiết bị khác nhau:
        \begin{itemize}
            \item Mobile app tương thích với hệ điều hành \textbf{Android 10 trở lên} và \textbf{iOS 14 trở lên}.
            \item Web app tương thích với các trình duyệt hiện đại như \textbf{Google Chrome}, \textbf{Mozilla Firefox}, \textbf{Microsoft Edge}, \textbf{Apple Safari} (đối với 2 phiên bản gần nhất).
        \end{itemize}
        \item Hệ thống tích hợp API với dịch vụ bên ngoài: social login (Google, Facebook), thanh toán (nếu có Premium), API âm nhạc (metadata, lyric).
    \end{itemize}

    % Tuân thủ pháp lý (Compliance)
    \item \textbf{Tuân thủ pháp lý (Compliance):} Hệ thống cần tuân thủ các quy định về bản quyền âm nhạc, bảo mật dữ liệu và quyền riêng tư.
    \begin{itemize}
        \item Hệ thống phải tuân thủ \textbf{Luật An ninh mạng Việt Nam} liên quan đến thu thập, xử lý, lưu trữ và bảo vệ thông tin cá nhân trong lãnh thổ Việt Nam.
        \item Nếu mở rộng quốc tế, phải đáp ứng chuẩn \textbf{GDPR (General Data Protection Regulation)} về bảo mật dữ liệu và quyền riêng tư.
        \item Đảm bảo bản quyền âm nhạc: chỉ lưu trữ, phát nhạc khi có giấy phép; có cơ chế gỡ bỏ nội dung vi phạm bản quyền theo yêu cầu trong vòng \textbf{48 giờ}.
    \end{itemize}

    % Khả năng kiểm toán (Audit)
    \item \textbf{Khả năng kiểm toán (Audit):}
    \begin{itemize}
        \item Hệ thống phải ghi lại toàn bộ hoạt động (logs) của người dùng và quản trị viên.
        \item Dữ liệu log cần được sao lưu định kỳ (ví dụ: ngày 28 hàng tháng).
        \item Log chỉ được đọc (read-only), không thể chỉnh sửa hay xóa qua giao diện người dùng.
    \end{itemize}

    % Hiệu quả sử dụng tài nguyên (Efficiency)
    \item \textbf{Hiệu quả sử dụng tài nguyên (Efficiency):}
    \begin{itemize}
        \item Hệ thống phải phục vụ ổn định với lưu lượng tối đa khoảng \textbf{5000 người dùng hoạt động đồng thời}.
        \item Tối ưu hóa tài nguyên hệ thống (RAM, CPU) để tránh lãng phí; sử dụng \textbf{adaptive bitrate (ABR)} để giảm băng thông khi mạng yếu.
        \item Ứng dụng mobile tối ưu pin: giảm bitrate khi màn hình tắt (audio-only mode).
    \end{itemize}

    % Khả năng bảo trì (Maintainability)
    \item \textbf{Khả năng bảo trì (Maintainability):}
    \begin{itemize}
        \item Có \textbf{tài liệu kỹ thuật} đầy đủ, rõ ràng (hướng dẫn cài đặt, vận hành, xử lý sự cố).
        \item Các lỗi phát sinh được ghi log và gửi thông báo cho kỹ thuật viên \textbf{trong vòng 5 phút} (qua email/sms).
        \item Hệ thống cần được kiểm tra và bảo trì định kỳ \textbf{tối đa 3 tháng/lần}.
        \item Kiến trúc phần mềm module hóa, hỗ trợ CI/CD, dễ dàng thêm mới tính năng như gợi ý nhạc AI, podcast,...
    \end{itemize}
\end{enumerate}

% Quản lý tài khoản - Người dùng
\begin{table}[h!]
\centering
\renewcommand{\arraystretch}{1.3} % tăng khoảng cách dòng trong bảng
\begin{tabularx}{\textwidth}{|l|X|}
\hline
\textbf{Use case name} & Quản lý tài khoản \\ \hline
\textbf{Created by}    & Vinh \\ \hline
\textbf{Actors}        & Người Dùng \\ \hline
\textbf{Description}   & Nơi để Người Dùng có thể quản lý tài khoản của mình (cập nhật thông tin cá nhân, đổi mật khẩu, đăng nhập, đăng ký, đăng xuất). \\ \hline
\textbf{Trigger}       & Người Dùng chọn chức năng “Quản lý tài khoản” trên giao diện hệ thống sau khi đăng nhập thành công. \\ \hline
\textbf{Pre-Condition(s)} 
& 1. Thiết bị của Người Dùng phải được kết nối internet. \newline
  2. Hệ thống hoạt động bình thường. \\ \hline
\textbf{Post-Condition(s)} 
& Thông tin tài khoản được cập nhật và lưu trữ an toàn trong hệ thống. \\ \hline
\textbf{Normal Flow}   
& 1. Người Dùng chọn vào phần “Quản lý tài khoản” trong giao diện hệ thống. \newline
  2. Hệ thống hiển thị thông tin tài khoản hiện tại. \newline
  3. Người Dùng chọn hành động: chỉnh sửa thông tin cá nhân / đổi mật khẩu / thiết lập bảo mật. \newline
  4. Người Dùng nhập thông tin mới hoặc thay đổi cần thiết. \newline
  5. Người Dùng xác nhận và lưu thay đổi. \newline
  6. Hệ thống cập nhật dữ liệu và thông báo thành công. \\ \hline
\textbf{Exception Flow} 
& 3a. Nếu hệ thống không tải được thông tin tài khoản thì hiển thị lỗi “Không thể tải dữ liệu”. \newline
  4a. Nếu thông tin nhập sai định dạng (ví dụ email không hợp lệ, mật khẩu quá ngắn), hệ thống thông báo lỗi và yêu cầu nhập lại. \newline
  5a. Nếu kết nối internet bị gián đoạn trong khi lưu, hệ thống hiển thị thông báo thất bại và yêu cầu thử lại. \\ \hline
\textbf{Alternative Flow} 
& 3b. Người Dùng có thể hủy thao tác và quay lại trang chính mà không thay đổi gì. \newline
  4b. Người Dùng có thể bật/tắt các tính năng nâng cao như xác thực 2 lớp. \newline
  5b. Người Dùng có thể tải xuống bản sao dữ liệu tài khoản của mình để lưu trữ. \\ \hline
\end{tabularx}
\caption{Mô tả usecase Quản lý tài khoản}
\end{table}

% Tương tác - Người dùng
\begin{table}[h!]
\centering
\renewcommand{\arraystretch}{1.3} % tăng khoảng cách dòng trong bảng
\begin{tabularx}{\textwidth}{|l|X|}
\hline
\textbf{Use case name} & Tương tác \\ \hline
\textbf{Created by}    & Vinh \\ \hline
\textbf{Actors}        & Người Dùng \\ \hline
\textbf{Description}   & Người Dùng có thể thực hiện các hành động tương tác trong hệ thống như: thích (like) bài hát, bình luận bài hát, theo dõi nghệ sĩ, chia sẻ bài hát/nghệ sĩ/playlist với bạn bè hoặc lên mạng xã hội. \\ \hline
\textbf{Trigger}       & Người Dùng chọn một bài hát, nghệ sĩ hoặc playlist bất kỳ trong hệ thống và mở giao diện chi tiết. \\ \hline
\textbf{Pre-Condition(s)} 
& 1. Thiết bị của Người Dùng phải được kết nối internet. \newline
  2. Người Dùng đã đăng nhập thành công. \\ \hline
\textbf{Post-Condition(s)} 
& Hệ thống lưu lại thông tin tương tác (like, bình luận, theo dõi, chia sẻ) và cập nhật dữ liệu liên quan. \\ \hline
\textbf{Normal Flow}   
& 1. Người Dùng chọn một bài hát/nghệ sĩ/playlist. \newline
  2. Hệ thống hiển thị chi tiết bài hát/nghệ sĩ/playlist. \newline
  3. Người Dùng chọn hành động muốn thực hiện: \newline
  \cach 3a. Nhấn nút “Thích” (Like) để thêm vào danh sách yêu thích. \newline
  \cach 3b. Viết và gửi bình luận cho bài hát. \newline
  \cach 3c. Nhấn “Theo dõi” để theo dõi nghệ sĩ. \newline
  \cach 3d. Chọn “Chia sẻ” và lựa chọn kênh chia sẻ (bạn bè trong ứng \cach dụng, mạng xã hội). \newline
  4. Hệ thống xác nhận và cập nhật thông tin tương tác. \\ \hline
\textbf{Exception Flow} 
& 2a. Nếu bài hát/nghệ sĩ/playlist không tồn tại hoặc bị xóa $\rightarrow$ hiển thị thông báo lỗi. \newline
  3b1. Nếu bình luận chứa nội dung vi phạm chính sách $\rightarrow$ hệ thống từ chối đăng và hiển thị thông báo. \newline
  3d1. Nếu chia sẻ thất bại do mất kết nối internet $\rightarrow$ hiển thị thông báo “Chia sẻ không thành công, vui lòng thử lại”. \\ \hline
\textbf{Alternative Flow} 
& 3a1. Người Dùng có thể “Bỏ thích” (Unlike) nếu trước đó đã thích. \newline
  3b2. Người Dùng có thể chỉnh sửa hoặc xóa bình luận đã đăng. \newline
  3c1. Người Dùng có thể hủy theo dõi nghệ sĩ. \newline
  3d2. Người Dùng có thể sao chép đường dẫn (link) thay vì chia sẻ trực tiếp. \\ \hline
\end{tabularx}
\caption{Mô tả usecase Tương tác}
\end{table}

% Nghe nhạc - Người dùng
\begin{table}[h!]
\centering
\renewcommand{\arraystretch}{1.3} % tăng khoảng cách dòng trong bảng
\begin{tabularx}{\textwidth}{|l|X|}
\hline
\textbf{Use case name} & Nghe nhạc \\ \hline
\textbf{Created by}    & Vinh \\ \hline
\textbf{Actors}        & Người Dùng \\ \hline
\textbf{Description}   & Người Dùng có thể phát và điều khiển nhạc với nhiều tính năng nâng cao: phát/tạm dừng, tiếp tục, chuyển tiếp/lùi lại, bật chế độ nghe ngẫu nhiên, nghe lặp lại (một bài hoặc toàn bộ danh sách), thêm bài vào hàng chờ, cài giờ tắt nhạc (sleep timer), và xem lời bài hát. \\ \hline
\textbf{Trigger}       & Người Dùng chọn một bài hát/playlist/album và nhấn nút “Phát nhạc”. \\ \hline
\textbf{Pre-Condition(s)} 
& 1. Thiết bị của Người Dùng phải được kết nối internet. \newline
  2. Ứng dụng được cấp quyền truy cập âm thanh. \\ \hline
\textbf{Post-Condition(s)} 
& Hệ thống phát nhạc theo thao tác của Người Dùng, cập nhật trạng thái trình phát (player state) và lưu lại lịch sử nghe nhạc. \\ \hline
\textbf{Normal Flow}   
& 1. Người Dùng chọn một bài hát/playlist/album. \newline
  2. Hệ thống tải dữ liệu nhạc và bắt đầu phát. \newline
  3. Người Dùng có thể thực hiện các thao tác điều khiển: \newline
  \cach 3a. Nhấn “Tạm dừng” (Pause) hoặc “Tiếp tục” (Play). \newline
  \cach 3b. Nhấn “Chuyển tiếp” (Next) hoặc “Lùi lại” (Previous). \newline
  \cach 3c. Bật chế độ “Nghe ngẫu nhiên” (Shuffle). \newline
  \cach 3d. Chọn chế độ “Nghe lặp lại” (Repeat one / Repeat all). \newline
  \cach 3e. Thêm bài hát vào hàng chờ phát (Queue). \newline
  \cach 3f. Cài đặt “Giờ tắt nhạc” (Sleep timer). \newline
  \cach 3g. Xem lời bài hát (Lyrics) nếu có sẵn. \newline
  4. Hệ thống phản hồi ngay lập tức và cập nhật trình phát nhạc. \\ \hline
\textbf{Exception Flow} 
& 2a. Nếu bài hát không thể phát (do bản quyền hoặc lỗi file) $\rightarrow$  hiển thị thông báo. \newline
  3f1. Nếu Người Dùng cài giờ tắt nhạc nhưng app bị thoát trước thời điểm đó $\rightarrow$ hệ thống không thể tắt nhạc đúng hẹn. \newline
  3g1. Nếu bài hát không có lời (lyrics) trong cơ sở dữ liệu $\rightarrow$ hiển thị “Chưa có lời bài hát”. \\ \hline
\textbf{Alternative Flow} 
& 3a1. Người Dùng có thể sử dụng tai nghe hoặc thiết bị ngoài để điều khiển (nút Play/Pause). \newline
  3e1. Người Dùng có thể sắp xếp lại thứ tự bài hát trong hàng chờ. \newline
  3f2. Người Dùng có thể hủy hoặc thay đổi thời gian sleep timer. \\ \hline
\end{tabularx}
\caption{Mô tả usecase Nghe nhạc}
\end{table}

% Khám phá và tìm kiếm - Người dùng
\begin{table}[h!]
\centering
\renewcommand{\arraystretch}{1.3} % tăng khoảng cách dòng trong bảng
\begin{tabularx}{\textwidth}{|l|X|}
\hline
\textbf{Use case name} & Khám phá và tìm kiếm \\ \hline
\textbf{Created by}    & Vinh \\ \hline
\textbf{Actors}        & Người Dùng \\ \hline
\textbf{Description}   & Người Dùng có thể tìm kiếm bài hát, nghệ sĩ, album, playlist và khám phá nhạc mới thông qua bảng xếp hạng, xu hướng (trending) và gợi ý cá nhân hóa. \\ \hline
\textbf{Trigger}       & Người Dùng mở thanh tìm kiếm hoặc tab “Khám phá” trong ứng dụng. \\ \hline
\textbf{Pre-Condition(s)} 
& Thiết bị của Người Dùng phải được kết nối internet. \\ \hline
\textbf{Post-Condition(s)} 
& Hệ thống trả về kết quả tìm kiếm hoặc danh sách nhạc khám phá phù hợp, cho phép Người Dùng chọn và nghe nhạc ngay. \\ \hline
\textbf{Normal Flow}   
& 1. Người Dùng mở tính năng Tìm kiếm/Khám phá. \newline
  2. Người Dùng có thể thực hiện: \newline
  \cach 2a. Nhập từ khóa để tìm bài hát, nghệ sĩ, album, playlist. \newline
  \cach 2b. Xem danh sách Trending hoặc Top Chart. \newline
  \cach 2c. Nhận gợi ý nhạc cá nhân hóa dựa trên lịch sử nghe, lượt thích và nghệ sĩ theo dõi. \newline
  3. Hệ thống hiển thị danh sách kết quả. \newline
  4. Người Dùng chọn nội dung mong muốn (phát nhạc, xem chi tiết nghệ sĩ/album/playlist). \\ \hline
\textbf{Exception Flow} 
& 2a1. Nếu không tìm thấy kết quả $\rightarrow$ hiển thị thông báo “Không tìm thấy nội dung phù hợp”. \newline
  2b1. Nếu hệ thống chưa có dữ liệu trending/top chart $\rightarrow$ hiển thị “Dữ liệu đang cập nhật”. \newline
  2c1. Nếu Người Dùng mới/ẩn danh (không đăng nhập) chưa có lịch sử nghe $\rightarrow$ gợi ý nhạc phổ biến mặc định. \\ \hline
\textbf{Alternative Flow} 
& 2a2. Người Dùng có thể dùng bộ lọc nâng cao (theo thể loại, thời lượng, năm phát hành…). \newline
  2b2. Người Dùng có thể chọn xem bảng xếp hạng theo từng khu vực/quốc gia. \newline
  2c2. Người Dùng có thể cập nhật gợi ý bằng cách thay đổi sở thích cá nhân (genres, mood…). \\ \hline
\end{tabularx}
\caption{Mô tả usecase Khám phá và tìm kiếm}
\end{table}
\end{document}